\documentclass[11pt,a4paper]{article}

\usepackage[T1]{fontenc}
\usepackage[utf8]{inputenc}
\usepackage{lmodern}
\usepackage{microtype}
\usepackage{amsmath,amssymb,amsthm,mathtools,bm}
\usepackage{graphicx}
\usepackage{booktabs}
\usepackage{enumitem}
\usepackage{xcolor}
\usepackage{geometry}
\usepackage{hyperref}
\usepackage{longtable}
\usepackage{array}

\geometry{margin=2.35cm}
\graphicspath{{../figures/}}
\hypersetup{colorlinks=true,linkcolor=blue!50!black,urlcolor=blue!50!black}
\setlength{\parindent}{0pt}
\setlength{\parskip}{0.42em}

\definecolor{accent}{RGB}{22,63,117}

\newcommand{\R}{\mathbb{R}}
\newcommand{\E}{\mathbb{E}}
\newcommand{\Var}{\mathrm{Var}}
\newcommand{\Cov}{\mathrm{Cov}}
\newcommand{\tr}{\mathrm{tr}}
\newcommand{\Dt}{\Delta_t}
\newcommand{\Ht}{\mathcal{H}_t}
\newcommand{\pt}{p_t}
\newcommand{\PhiN}{\Phi}
\newcommand{\phiv}{\phi_v}
\newcommand{\One}{\mathbf 1}

\newtheorem{proposition}{Proposition}
\newtheorem{remark}{Remark}

\title{\textbf{\color{accent}Derivation compendium and audit notes}\\[2mm]
\Large corrected companion for the Laplace-focused presentation}
\author{}
\date{March 2026}

\begin{document}
\maketitle
\vspace{-1.0em}

\begin{center}
\begin{minipage}{0.97\textwidth}
\small
This companion records the main corrections to the original draft and proves every formula used in the short presentation note. The emphasis is on explicit algebra: products of Gaussians are completed by hand, matrix identities are shown line by line, and every derivative used in the $K=1$ Laplace benchmark is written out.
\end{minipage}
\end{center}

\tableofcontents
\bigskip

\section{Scope and main corrections}

The clean benchmark retained in the final short note is the two-frame Laplace AR(1) model
\begin{equation}
A_0 \sim \mathcal N(\mu_0,\sigma_0^2),
\qquad
A_1 = \alpha A_0 + \varepsilon,
\qquad
\varepsilon \sim \mathrm{Laplace}(0,b),
\qquad |\alpha|<1,
\label{eq:setup-clean}
\end{equation}
followed by OU noising
\begin{equation}
X_k = e^{-t} A_k + \sqrt{\Dt} \, Z_k,
\qquad
\Dt = 1-e^{-2t},
\qquad
Z_0,Z_1 \stackrel{\mathrm{iid}}{\sim} \mathcal N(0,1).
\label{eq:setup-ou}
\end{equation}

The original draft contained several mathematical issues. The final package corrects them as follows.
\begin{enumerate}[leftmargin=1.45em,itemsep=0.25em]
\item The notation $a_0 \sim \mathcal N(\mu_0,\Sigma_0)$ was inconsistent: for the scalar initial variable one needs $a_0 \sim \mathcal N(\mu_0,\sigma_0^2)$, while $\Sigma_0$ denotes the covariance of the full trajectory vector.
\item The displayed finite-chain Gaussian AR(1) precision matrix was not exact at the boundaries. The first and last diagonal entries differ from the interior entries; the correct matrix is derived in Section \ref{sec:gaussian-precision}.
\item The locality diagnostics based on the full Frobenius norm of $Q_t$ were misleading at large time, because $Q_t \to I$. The corrected diagnostics normalize only off-diagonal mass; see Section \ref{sec:locality-diagnostics}.
\item The deterministic control was mathematically correct but conceptually overinterpreted. Its dense noised precision comes from a rank-one propagation mechanism, not from the same Markov fill-in mechanism as the stochastic AR(1) chain; see Section \ref{sec:det-control}.
\item The old normal-Laplace convolution formula used in the code was wrong. The corrected formula is derived in Section \ref{sec:normal-laplace}.
\item The old $K=2$ Laplace ``closed-form'' claim is not retained. What is justified is an exact nested integral representation, not an elementary closed form comparable to the $K=1$ formula; see Section \ref{sec:k2-statement}.
\end{enumerate}

The rest of the document proves the corrected statements explicitly.

\section{Gaussian and deterministic audit}

\subsection{Exact finite-chain Gaussian AR(1) precision}
\label{sec:gaussian-precision}

Consider the Gaussian AR(1) chain
\begin{equation}
A_{k+1} = \alpha A_k + \eta_k,
\qquad
\eta_k \sim \mathcal N(0,\sigma_\eta^2),
\qquad
k=0,\dots,K-1,
\label{eq:gaussian-ar1}
\end{equation}
with
\begin{equation}
A_0 \sim \mathcal N(\mu_0,\sigma_0^2).
\label{eq:a0-gaussian}
\end{equation}
Its joint density is
\begin{equation}
p(a_0,\dots,a_K)
\propto
\exp\!\left(
-\frac{(a_0-\mu_0)^2}{2\sigma_0^2}
-\sum_{k=0}^{K-1}\frac{(a_{k+1}-\alpha a_k)^2}{2\sigma_\eta^2}
\right).
\label{eq:joint-gaussian-chain}
\end{equation}

We expand the quadratic form in \eqref{eq:joint-gaussian-chain} exactly.

\paragraph{Step 1: expand the initial term.}
\begin{equation}
\frac{(a_0-\mu_0)^2}{2\sigma_0^2}
=
\frac{a_0^2}{2\sigma_0^2}
-\frac{\mu_0}{\sigma_0^2} a_0
+\frac{\mu_0^2}{2\sigma_0^2}.
\label{eq:initial-expand}
\end{equation}

\paragraph{Step 2: expand each innovation term.}
For each $k$,
\begin{equation}
\frac{(a_{k+1}-\alpha a_k)^2}{2\sigma_\eta^2}
=
\frac{a_{k+1}^2}{2\sigma_\eta^2}
-\frac{\alpha}{\sigma_\eta^2} a_k a_{k+1}
+\frac{\alpha^2 a_k^2}{2\sigma_\eta^2}.
\label{eq:innovation-expand}
\end{equation}

\paragraph{Step 3: collect the coefficients of each $a_j^2$.}
\begin{itemize}[leftmargin=1.4em]
\item The variable $a_0$ appears in \eqref{eq:initial-expand} and in the $k=0$ instance of \eqref{eq:innovation-expand}. Therefore its quadratic coefficient is
\[
\frac{1}{2\sigma_0^2} + \frac{\alpha^2}{2\sigma_\eta^2}.
\]
\item An interior variable $a_j$, with $1\le j\le K-1$, appears
\begin{itemize}[leftmargin=1.2em]
\item as $a_{k+1}$ when $k=j-1$, giving $a_j^2/(2\sigma_\eta^2)$,
\item as $a_k$ when $k=j$, giving $\alpha^2 a_j^2/(2\sigma_\eta^2)$.
\end{itemize}
Therefore its quadratic coefficient is
\[
\frac{1+\alpha^2}{2\sigma_\eta^2}.
\]
\item The last variable $a_K$ appears only as $a_{k+1}$ when $k=K-1$. Therefore its quadratic coefficient is
\[
\frac{1}{2\sigma_\eta^2}.
\]
\end{itemize}

\paragraph{Step 4: collect cross terms.}
From \eqref{eq:innovation-expand}, the only cross term involving $a_k$ and $a_{k+1}$ is
\begin{equation}
-\frac{\alpha}{\sigma_\eta^2} a_k a_{k+1}.
\label{eq:cross-term}
\end{equation}
No other cross terms appear. Hence the precision matrix is tridiagonal.

Putting these pieces together,
\begin{equation}
p(a) \propto \exp\!\left(-\frac12 a^\top Q_0 a + h^\top a\right),
\label{eq:gauss-quad-form}
\end{equation}
with precision
\begin{equation}
Q_0 =
\begin{pmatrix}
\sigma_0^{-2}+\alpha^2 \sigma_\eta^{-2} & -\alpha \sigma_\eta^{-2} & & & \\
-\alpha \sigma_\eta^{-2} & (1+\alpha^2)\sigma_\eta^{-2} & -\alpha \sigma_\eta^{-2} & & \\
& \ddots & \ddots & \ddots & \\
& & -\alpha \sigma_\eta^{-2} & (1+\alpha^2)\sigma_\eta^{-2} & -\alpha \sigma_\eta^{-2}\\
& & & -\alpha \sigma_\eta^{-2} & \sigma_\eta^{-2}
\end{pmatrix}.
\label{eq:Q0-general}
\end{equation}

\paragraph{Stationary specialization.}
If $A_0$ is chosen with stationary variance
\begin{equation}
\sigma_0^2 = \frac{\sigma_\eta^2}{1-\alpha^2},
\label{eq:stationary-variance}
\end{equation}
then
\[
\sigma_0^{-2}+\alpha^2 \sigma_\eta^{-2}
=
\frac{1-\alpha^2}{\sigma_\eta^2} + \frac{\alpha^2}{\sigma_\eta^2}
=
\frac{1}{\sigma_\eta^2}.
\]
Hence the first and last diagonal entries coincide, and
\begin{equation}
Q_0 =
\frac{1}{\sigma_\eta^2}
\begin{pmatrix}
1 & -\alpha & & & \\
-\alpha & 1+\alpha^2 & -\alpha & & \\
& \ddots & \ddots & \ddots & \\
& & -\alpha & 1+\alpha^2 & -\alpha\\
& & & -\alpha & 1
\end{pmatrix}.
\label{eq:Q0-stationary}
\end{equation}
This is the correct finite-chain formula. The original short draft displayed the interior pattern on every diagonal entry, which is not exact at the boundaries.

\subsection{OU noising and the exact Gaussian precision identities}
\label{sec:ou-identities}

Let $A\in\R^n$ be any random vector with mean $\mu_0$ and covariance $\Sigma_0$. Define
\begin{equation}
X_t = e^{-t} A + \sqrt{\Dt} \, Z,
\qquad
Z \sim \mathcal N(0,I_n),
\qquad Z \perp A.
\label{eq:generic-ou}
\end{equation}
Then
\begin{align}
\E[X_t]
&= e^{-t}\E[A] + \sqrt{\Dt}\,\E[Z]
= e^{-t}\mu_0,
\label{eq:ou-mean-derivation}\\[0.4em]
\Cov(X_t)
&=
\Cov\bigl(e^{-t}A + \sqrt{\Dt} Z\bigr)
\notag\\
&=
 e^{-2t}\Cov(A)
 + \Dt \Cov(Z)
 + e^{-t}\sqrt{\Dt}\,\Cov(A,Z)
 + e^{-t}\sqrt{\Dt}\,\Cov(Z,A).
\label{eq:ou-cov-expand}
\end{align}
Because $A$ and $Z$ are independent, $\Cov(A,Z)=\Cov(Z,A)=0$, and because $\Cov(Z)=I_n$, \eqref{eq:ou-cov-expand} becomes
\begin{equation}
\Sigma_t := \Cov(X_t) = e^{-2t} \Sigma_0 + \Dt I_n.
\label{eq:sigma-t}
\end{equation}

\paragraph{SNR form of the precision.}
Let
\begin{equation}
\gamma_t := \frac{e^{-2t}}{\Dt}.
\label{eq:gamma-t}
\end{equation}
Then
\[
\Sigma_t = \Dt \bigl(I_n + \gamma_t \Sigma_0\bigr),
\]
so
\begin{equation}
Q_t := \Sigma_t^{-1}
=
\Dt^{-1} \bigl(I_n + \gamma_t \Sigma_0\bigr)^{-1}.
\label{eq:Qt-snr}
\end{equation}

\paragraph{Factorization through $Q_0=\Sigma_0^{-1}$.}
Since $\Dt = 1-e^{-2t}$, one has
\begin{equation}
e^{2t}\Dt = e^{2t}-1 =: c_t.
\label{eq:ct-def}
\end{equation}
Rewrite \eqref{eq:sigma-t} as
\begin{equation}
\Sigma_t = e^{-2t}\bigl(\Sigma_0 + c_t I_n\bigr).
\label{eq:sigma-factor}
\end{equation}
Invert \eqref{eq:sigma-factor}:
\begin{equation}
Q_t = e^{2t} (\Sigma_0 + c_t I_n)^{-1}.
\label{eq:Qt-first-factor}
\end{equation}
Now use $\Sigma_0 = Q_0^{-1}$:
\[
\Sigma_0 + c_t I_n
=
Q_0^{-1} + c_t I_n
=
Q_0^{-1}(I_n + c_t Q_0).
\]
Hence
\[
(\Sigma_0 + c_t I_n)^{-1}
=
(I_n + c_t Q_0)^{-1} Q_0.
\]
Substitute into \eqref{eq:Qt-first-factor}:
\begin{equation}
Q_t = e^{2t} (I_n + c_t Q_0)^{-1} Q_0.
\label{eq:Qt-factor-right}
\end{equation}
Because $Q_0$ commutes with every polynomial in $Q_0$, it also commutes with $(I_n+c_tQ_0)^{-1}$. Therefore
\begin{equation}
Q_t = e^{2t} Q_0 (I_n + c_t Q_0)^{-1}.
\label{eq:Qt-factor-left}
\end{equation}

\paragraph{Neumann expansion.}
If $c_t \|Q_0\|<1$ in any submultiplicative operator norm, then
\begin{equation}
(I_n + c_t Q_0)^{-1}
=
\sum_{m=0}^{\infty} (-c_t)^m Q_0^m.
\label{eq:neumann}
\end{equation}
Insert this into \eqref{eq:Qt-factor-left}:
\begin{equation}
Q_t = e^{2t} \sum_{m=0}^{\infty} (-c_t)^m Q_0^{m+1}.
\label{eq:Qt-neumann}
\end{equation}

\paragraph{Why distance-$d$ couplings appear at order $t^{d-1}$.}
If $Q_0$ is tridiagonal, then $Q_0$ has bandwidth $1$: $(Q_0)_{ij}=0$ whenever $|i-j|>1$.

If $B$ has bandwidth $r$ and $C$ has bandwidth $s$, then $BC$ has bandwidth at most $r+s$. Indeed,
\[
(BC)_{ij} = \sum_k B_{ik} C_{kj}.
\]
For a term in the sum to be nonzero, one must have both $|i-k|\le r$ and $|k-j|\le s$. Therefore $|i-j|\le r+s$. Hence $Q_0^{m+1}$ has bandwidth at most $m+1$.

In particular, an entry with $|i-j|=d$ cannot appear before the term $Q_0^d$, which corresponds to $m=d-1$ in \eqref{eq:Qt-neumann}. For small $t$,
\[
c_t = e^{2t}-1 = 2t + O(t^2),
\]
so the first nonzero contribution to a distance-$d$ entry is of order $t^{d-1}$.

\paragraph{Large-time expansion.}
Let $\varepsilon = e^{-2t}$. Then $\Dt = 1-\varepsilon$, and \eqref{eq:sigma-t} becomes
\[
\Sigma_t = I_n + \varepsilon (\Sigma_0-I_n).
\]
For small $\varepsilon$,
\[
(I_n + \varepsilon M)^{-1} = I_n - \varepsilon M + O(\varepsilon^2).
\]
Apply this with $M=\Sigma_0-I_n$:
\begin{equation}
Q_t = I_n - e^{-2t}(\Sigma_0-I_n) + O(e^{-4t}).
\label{eq:Qt-large-time}
\end{equation}

\subsection{Gaussian score and posterior mean}

If $X_t \sim \mathcal N(\mu_t,\Sigma_t)$ with precision $Q_t=\Sigma_t^{-1}$, then
\begin{equation}
p_t^{\mathrm G}(x)
=
\frac{1}{(2\pi)^{n/2} (\det \Sigma_t)^{1/2}}
\exp\!\left(-\frac12 (x-\mu_t)^\top Q_t (x-\mu_t)\right).
\label{eq:gaussian-density-general}
\end{equation}
Take logarithms:
\[
\log p_t^{\mathrm G}(x)
=
C_t - \frac12 (x-\mu_t)^\top Q_t (x-\mu_t),
\]
where $C_t$ does not depend on $x$. Expand the quadratic term:
\[
(x-\mu_t)^\top Q_t (x-\mu_t)
=
 x^\top Q_t x - 2 \mu_t^\top Q_t x + \mu_t^\top Q_t \mu_t.
\]
Therefore
\[
\log p_t^{\mathrm G}(x)
=
C_t' - \frac12 x^\top Q_t x + \mu_t^\top Q_t x.
\]
Differentiate with respect to $x$:
\begin{equation}
\nabla \log p_t^{\mathrm G}(x) = -Q_t(x-\mu_t),
\qquad
-\nabla^2 \log p_t^{\mathrm G}(x) \equiv Q_t.
\label{eq:gaussian-score-hessian}
\end{equation}

The posterior mean of the clean state given $X_t=x$ is the standard Gaussian regression formula. Since
\[
\Cov(A,X_t) = \Cov\bigl(A,e^{-t}A+\sqrt{\Dt}Z\bigr) = e^{-t}\Sigma_0,
\]
and $\Cov(X_t)=\Sigma_t$, one gets
\begin{equation}
\E[A\mid X_t=x]
=
\mu_A + e^{-t}\Sigma_0 Q_t (x-\mu_t).
\label{eq:posterior-mean-gaussian}
\end{equation}

\subsection{Corrected locality diagnostics}
\label{sec:locality-diagnostics}

The original short draft normalized band energies by $\|Q_t\|_F^2$ including the diagonal. This becomes misleading at large $t$, because $Q_t \to I$, so the diagonal dominates automatically.

A more meaningful off-diagonal normalization is
\begin{equation}
E_d(t) := \sum_{|i-j|=d} Q_t[i,j]^2,
\qquad
E_{\mathrm{off}}(t) := \sum_{d\ge 1} E_d(t),
\qquad
p_d(t) := \frac{E_d(t)}{E_{\mathrm{off}}(t)}.
\label{eq:offdiag-diagnostics}
\end{equation}
From these one can define
\begin{equation}
b_{0.95}(t) := \min\left\{r:\sum_{d=1}^r p_d(t)\ge 0.95\right\},
\qquad
\bar\xi(t) := \sum_{d\ge 1} d\, p_d(t).
\label{eq:bandwidth-range}
\end{equation}
These observables measure the shape of the off-diagonal coupling profile rather than the trivial approach to the identity.

\subsection{Deterministic control}
\label{sec:det-control}

The deterministic control is
\begin{equation}
A_k = \alpha^k A_0,
\qquad
A_0 \sim \mathcal N(0,\sigma_0^2),
\qquad
v := (1,\alpha,\dots,\alpha^K)^\top.
\label{eq:det-control}
\end{equation}
Then the clean covariance is rank one:
\begin{equation}
\Sigma_0^{\mathrm{det}} = \sigma_0^2 v v^\top.
\label{eq:det-cov-clean}
\end{equation}
After OU noising,
\begin{equation}
\Sigma_t^{\mathrm{det}} = \Dt I + \beta_t v v^\top,
\qquad
\beta_t := e^{-2t} \sigma_0^2.
\label{eq:det-cov-noisy}
\end{equation}
We invert \eqref{eq:det-cov-noisy} by Sherman-Morrison. Write
\[
A = \Dt I,
\qquad
u = \sqrt{\beta_t}\,v,
\qquad
\Sigma_t^{\mathrm{det}} = A + u u^\top.
\]
The Sherman-Morrison formula gives
\begin{equation}
(A+u u^\top)^{-1}
=
A^{-1} - \frac{A^{-1} u u^\top A^{-1}}{1+u^\top A^{-1} u}.
\label{eq:sherman-morrison}
\end{equation}
Because $A^{-1}=\Dt^{-1} I$,
\[
A^{-1}u = \Dt^{-1} \sqrt{\beta_t}\, v,
\qquad
u^\top A^{-1}u = \frac{\beta_t}{\Dt} \|v\|^2.
\]
Also,
\[
A^{-1}u u^\top A^{-1}
=
\frac{\beta_t}{\Dt^2} v v^\top.
\]
Substitute into \eqref{eq:sherman-morrison}:
\[
(\Sigma_t^{\mathrm{det}})^{-1}
=
\frac{1}{\Dt} I - \frac{(\beta_t/\Dt^2) v v^\top}{1+(\beta_t/\Dt)\|v\|^2}.
\]
Multiply numerator and denominator of the fraction by $\Dt$:
\begin{equation}
Q_t^{\mathrm{det}}
=
\frac{1}{\Dt} I - \frac{\beta_t}{\Dt(\Dt+\beta_t\|v\|^2)} v v^\top.
\label{eq:det-precision}
\end{equation}
This matrix is dense for every $t>0$, because $v v^\top$ is dense. The corrected interpretation is:
\begin{quote}
The deterministic control shows that a dense noised precision can arise from a rank-one propagation mechanism. This is different from the specific Gaussian-Markov band fill-in mechanism of the stochastic AR(1) chain.
\end{quote}

\section[Laplace K=1: exact characteristic function]{Laplace $K=1$: exact characteristic function}
\label{sec:cf-laplace}

\subsection{Clean characteristic function}

Let $\xi=(\xi_0,\xi_1)\in\R^2$. By definition,
\[
\Phi_0(\xi_0,\xi_1) = \E\!\left[e^{i(\xi_0 A_0 + \xi_1 A_1)}\right].
\]
Since $A_1 = \alpha A_0 + \varepsilon$,
\[
\xi_0 A_0 + \xi_1 A_1 = (\xi_0+\alpha\xi_1) A_0 + \xi_1 \varepsilon.
\]
Because $A_0$ and $\varepsilon$ are independent,
\begin{equation}
\Phi_0(\xi_0,\xi_1)
=
\E\!\left[e^{i(\xi_0+\alpha\xi_1)A_0}\right]
\cdot
\E\!\left[e^{i\xi_1 \varepsilon}\right].
\label{eq:cf-factor-clean}
\end{equation}

\paragraph{Gaussian factor.}
If $Y\sim\mathcal N(\mu,\sigma^2)$, then
\[
\E[e^{isY}] = \exp\!\left(is\mu - \frac12 \sigma^2 s^2\right).
\]
Apply this with $Y=A_0$ and $s=\xi_0+\alpha\xi_1$:
\begin{equation}
\E\!\left[e^{i(\xi_0+\alpha\xi_1)A_0}\right]
=
\exp\!\left(
 i\mu_0(\xi_0+\alpha\xi_1)
 - \frac12 \sigma_0^2 (\xi_0+\alpha\xi_1)^2
\right).
\label{eq:cf-a0}
\end{equation}

\paragraph{Laplace factor.}
We compute the characteristic function of $\varepsilon\sim\mathrm{Laplace}(0,b)$:
\begin{align}
\E[e^{is\varepsilon}]
&=
\int_{-\infty}^{\infty} e^{isu} \frac{e^{-|u|/b}}{2b}\,du
\notag\\
&=
\frac{1}{2b}\int_{-\infty}^{0} e^{isu} e^{u/b}\,du
+
\frac{1}{2b}\int_{0}^{\infty} e^{isu} e^{-u/b}\,du.
\label{eq:laplace-cf-split}
\end{align}
The first integral is
\[
\int_{-\infty}^{0} e^{(is+1/b)u}\,du
=
\left[\frac{e^{(is+1/b)u}}{is+1/b}\right]_{u=-\infty}^{u=0}
=
\frac{1}{is+1/b},
\]
because $\Re(is+1/b)=1/b>0$, so the exponential vanishes at $-\infty$.

The second integral is
\[
\int_{0}^{\infty} e^{(is-1/b)u}\,du
=
\left[\frac{e^{(is-1/b)u}}{is-1/b}\right]_{u=0}^{u=\infty}
=
-\frac{1}{is-1/b}
=
\frac{1}{1/b-is}.
\]
Hence
\[
\E[e^{is\varepsilon}]
=
\frac{1}{2b}
\left(
\frac{1}{is+1/b}
+
\frac{1}{1/b-is}
\right).
\]
Put the two fractions over a common denominator:
\[
\frac{1}{is+1/b} + \frac{1}{1/b-is}
=
\frac{(1/b-is)+(is+1/b)}{(is+1/b)(1/b-is)}
=
\frac{2/b}{1/b^2+s^2}.
\]
Therefore
\begin{equation}
\E[e^{is\varepsilon}] = \frac{1}{1+b^2 s^2}.
\label{eq:laplace-cf}
\end{equation}

Combining \eqref{eq:cf-factor-clean}, \eqref{eq:cf-a0}, and \eqref{eq:laplace-cf},
\begin{equation}
\Phi_0(\xi_0,\xi_1)
=
\exp\!\left(
 i\mu_0(\xi_0+\alpha\xi_1)
 - \frac12 \sigma_0^2 (\xi_0+\alpha\xi_1)^2
\right)
\frac{1}{1+b^2 \xi_1^2}.
\label{eq:cf-clean-final}
\end{equation}

\subsection{Characteristic function after OU noising}

Conditionally on $A=(A_0,A_1)$, one has
\[
X_t = e^{-t} A + \sqrt{\Dt} Z.
\]
Hence
\begin{align}
\E[e^{i\xi^\top X_t}\mid A]
&=
\E[e^{i\xi^\top(e^{-t}A+\sqrt{\Dt}Z)}\mid A]
\notag\\
&=
 e^{i e^{-t}\xi^\top A}
\, \E[e^{i\sqrt{\Dt}\,\xi^\top Z}]
\notag\\
&=
 e^{i e^{-t}\xi^\top A}
\, e^{-\Dt \|\xi\|^2/2}.
\label{eq:cf-conditional-noise}
\end{align}
Taking expectation in $A$,
\begin{equation}
\Phi_t(\xi) = e^{-\Dt\|\xi\|^2/2} \, \Phi_0(e^{-t}\xi).
\label{eq:cf-ou-rule}
\end{equation}
Insert \eqref{eq:cf-clean-final}:
\begin{equation}
\Phi_t(\xi_0,\xi_1)
=
\exp\!\Bigl(
 i e^{-t}\mu_0(\xi_0+\alpha\xi_1)
 -\frac12 \Dt (\xi_0^2+\xi_1^2)
 -\frac12 e^{-2t}\sigma_0^2 (\xi_0+\alpha\xi_1)^2
\Bigr)
\frac{1}{1+b^2 e^{-2t} \xi_1^2}.
\label{eq:cf-final}
\end{equation}

A centered Gaussian characteristic function is always of the form $\exp(-\xi^\top \Sigma \xi /2)$, possibly multiplied by a pure phase $e^{i m^\top \xi}$ if the mean is not zero. The extra factor
\[
\frac{1}{1+b^2 e^{-2t} \xi_1^2}
\]
in \eqref{eq:cf-final} is rational, not exponential-quadratic. Therefore the diffused Laplace law is not Gaussian for any finite $t$.

\section{Derivation of the normal-Laplace convolution}
\label{sec:normal-laplace}

Fix $\beta>0$ and $\sigma^2>0$. Define
\begin{equation}
h(y;\beta,\sigma^2)
=
\int_{-\infty}^{\infty}
\frac{e^{-|u|/\beta}}{2\beta}
\cdot
\frac{1}{\sqrt{2\pi\sigma^2}}
\exp\!\left(-\frac{(y-u)^2}{2\sigma^2}\right)du.
\label{eq:h-start}
\end{equation}
This is the density of $L+G$ with $L\sim\mathrm{Laplace}(0,\beta)$ and $G\sim\mathcal N(0,\sigma^2)$ independent.

\subsection[Split the integral at u=0]{Split the integral at $u=0$}

Write
\begin{equation}
h(y;\beta,\sigma^2)
=
\frac{1}{2\beta\sqrt{2\pi\sigma^2}} \Bigl(I_+(y)+I_-(y)\Bigr),
\label{eq:h-split}
\end{equation}
where
\begin{align}
I_+(y)
&=
\int_0^{\infty}
\exp\!\left(-\frac{u}{\beta} - \frac{(y-u)^2}{2\sigma^2}\right)du,
\label{eq:Iplus-def}\\[0.3em]
I_-(y)
&=
\int_{-\infty}^{0}
\exp\!\left(\frac{u}{\beta} - \frac{(y-u)^2}{2\sigma^2}\right)du.
\label{eq:Iminus-def}
\end{align}

\subsection[Complete the square in I+]{Complete the square in $I_+$}

Start with the exponent:
\begin{align}
-\frac{u}{\beta} - \frac{(y-u)^2}{2\sigma^2}
&=
-\frac{u}{\beta} - \frac{u^2-2yu+y^2}{2\sigma^2}
\notag\\
&=
-\frac{1}{2\sigma^2}
\left(
 u^2 - 2yu + y^2 + \frac{2\sigma^2}{\beta}u
\right)
\notag\\
&=
-\frac{1}{2\sigma^2}
\left(
 u^2 - 2\left(y-\frac{\sigma^2}{\beta}\right)u + y^2
\right).
\label{eq:Iplus-square1}
\end{align}
Now insert and subtract $\bigl(y-\sigma^2/\beta\bigr)^2$:
\begin{align}
u^2 - 2\left(y-\frac{\sigma^2}{\beta}\right)u + y^2
&=
\left(u-\left(y-\frac{\sigma^2}{\beta}\right)\right)^2
+
y^2 - \left(y-\frac{\sigma^2}{\beta}\right)^2
\notag\\
&=
\left(u-\left(y-\frac{\sigma^2}{\beta}\right)\right)^2
+
\left(y^2 - y^2 + \frac{2y\sigma^2}{\beta} - \frac{\sigma^4}{\beta^2}\right)
\notag\\
&=
\left(u-\left(y-\frac{\sigma^2}{\beta}\right)\right)^2
+
\frac{2y\sigma^2}{\beta} - \frac{\sigma^4}{\beta^2}.
\label{eq:Iplus-square2}
\end{align}
Substitute \eqref{eq:Iplus-square2} into \eqref{eq:Iplus-square1}:
\begin{equation}
-\frac{u}{\beta} - \frac{(y-u)^2}{2\sigma^2}
=
-\frac{\left(u-(y-\sigma^2/\beta)\right)^2}{2\sigma^2}
-\frac{y}{\beta} + \frac{\sigma^2}{2\beta^2}.
\label{eq:Iplus-square3}
\end{equation}
Therefore
\begin{equation}
I_+(y)
=
 e^{-y/\beta+\sigma^2/(2\beta^2)}
\int_0^{\infty}
\exp\!\left(-\frac{(u-(y-\sigma^2/\beta))^2}{2\sigma^2}\right)du.
\label{eq:Iplus-after-square}
\end{equation}
Set
\[
z = \frac{u-(y-\sigma^2/\beta)}{\sigma},
\qquad du = \sigma dz.
\]
When $u=0$,
\[
z = -\frac{y}{\sigma} + \frac{\sigma}{\beta},
\]
and when $u\to\infty$, $z\to\infty$. Hence
\begin{align}
I_+(y)
&=
 e^{-y/\beta+\sigma^2/(2\beta^2)}
\sigma \int_{-y/\sigma+\sigma/\beta}^{\infty} e^{-z^2/2} dz
\notag\\
&=
 e^{-y/\beta+\sigma^2/(2\beta^2)}
\sigma\sqrt{2\pi}\,\PhiN\!\left(\frac{y}{\sigma}-\frac{\sigma}{\beta}\right).
\label{eq:Iplus-final}
\end{align}

\subsection[Complete the square in I-]{Complete the square in $I_-$}

Now start from
\begin{align}
\frac{u}{\beta} - \frac{(y-u)^2}{2\sigma^2}
&=
\frac{u}{\beta} - \frac{u^2-2yu+y^2}{2\sigma^2}
\notag\\
&=
-\frac{1}{2\sigma^2}
\left(
 u^2 - 2yu + y^2 - \frac{2\sigma^2}{\beta}u
\right)
\notag\\
&=
-\frac{1}{2\sigma^2}
\left(
 u^2 - 2\left(y+\frac{\sigma^2}{\beta}\right)u + y^2
\right).
\label{eq:Iminus-square1}
\end{align}
Insert and subtract $\bigl(y+\sigma^2/\beta\bigr)^2$:
\begin{align}
u^2 - 2\left(y+\frac{\sigma^2}{\beta}\right)u + y^2
&=
\left(u-\left(y+\frac{\sigma^2}{\beta}\right)\right)^2
+
y^2 - \left(y+\frac{\sigma^2}{\beta}\right)^2
\notag\\
&=
\left(u-\left(y+\frac{\sigma^2}{\beta}\right)\right)^2
-
\frac{2y\sigma^2}{\beta} - \frac{\sigma^4}{\beta^2}.
\label{eq:Iminus-square2}
\end{align}
Substitute into \eqref{eq:Iminus-square1}:
\begin{equation}
\frac{u}{\beta} - \frac{(y-u)^2}{2\sigma^2}
=
-\frac{\left(u-(y+\sigma^2/\beta)\right)^2}{2\sigma^2}
+\frac{y}{\beta} + \frac{\sigma^2}{2\beta^2}.
\label{eq:Iminus-square3}
\end{equation}
Hence
\begin{equation}
I_-(y)
=
 e^{y/\beta+\sigma^2/(2\beta^2)}
\int_{-\infty}^{0}
\exp\!\left(-\frac{(u-(y+\sigma^2/\beta))^2}{2\sigma^2}\right)du.
\label{eq:Iminus-after-square}
\end{equation}
Set
\[
z = \frac{u-(y+\sigma^2/\beta)}{\sigma},
\qquad du = \sigma dz.
\]
When $u=0$,
\[
z = -\frac{y}{\sigma} - \frac{\sigma}{\beta},
\]
and when $u\to-\infty$, $z\to-\infty$. Therefore
\begin{align}
I_-(y)
&=
 e^{y/\beta+\sigma^2/(2\beta^2)}
\sigma \int_{-\infty}^{-y/\sigma-\sigma/\beta} e^{-z^2/2} dz
\notag\\
&=
 e^{y/\beta+\sigma^2/(2\beta^2)}
\sigma\sqrt{2\pi}\,\PhiN\!\left(-\frac{y}{\sigma}-\frac{\sigma}{\beta}\right).
\label{eq:Iminus-final}
\end{align}

\subsection{Final formula and an equivalent erfc form}

Insert \eqref{eq:Iplus-final} and \eqref{eq:Iminus-final} into \eqref{eq:h-split}:
\begin{equation}
h(y;\beta,\sigma^2)
=
\frac{e^{\sigma^2/(2\beta^2)}}{2\beta}
\left[
 e^{-y/\beta}\PhiN\!\left(\frac{y}{\sigma}-\frac{\sigma}{\beta}\right)
 +
 e^{y/\beta}\PhiN\!\left(-\frac{y}{\sigma}-\frac{\sigma}{\beta}\right)
\right].
\label{eq:h-final}
\end{equation}
Using $\PhiN(z)=\frac12 \operatorname{erfc}(-z/\sqrt{2})$, one also gets
\begin{equation}
h(y;\beta,\sigma^2)
=
\frac{e^{\sigma^2/(2\beta^2)}}{4\beta}
\left[
 e^{-y/\beta}\operatorname{erfc}\!\left(\frac{\sigma^2/\beta-y}{\sqrt{2}\sigma}\right)
 +
 e^{y/\beta}\operatorname{erfc}\!\left(\frac{\sigma^2/\beta+y}{\sqrt{2}\sigma}\right)
\right].
\label{eq:h-erfc}
\end{equation}
This is the corrected formula used in the new code.

\subsection[Derivative formulas for h]{Derivative formulas for $h$}

Set
\[
u = \frac{y}{\sigma} - \frac{\sigma}{\beta},
\qquad
v = -\frac{y}{\sigma} - \frac{\sigma}{\beta},
\qquad
C = \frac{e^{\sigma^2/(2\beta^2)}}{2\beta}.
\]
Then \eqref{eq:h-final} reads
\[
h(y) = C \bigl(e^{-y/\beta}\PhiN(u) + e^{y/\beta}\PhiN(v)\bigr).
\]
Differentiate the first term:
\begin{align}
\frac{d}{dy}\left[e^{-y/\beta}\PhiN(u)\right]
&=
-\frac{1}{\beta} e^{-y/\beta}\PhiN(u)
+
 e^{-y/\beta} \phi(u) \frac{du}{dy}
\notag\\
&=
-\frac{1}{\beta} e^{-y/\beta}\PhiN(u)
+
\frac{1}{\sigma} e^{-y/\beta} \phi(u),
\label{eq:hprime-term1}
\end{align}
where $\phi(z)=\frac{1}{\sqrt{2\pi}} e^{-z^2/2}$.
Differentiate the second term:
\begin{align}
\frac{d}{dy}\left[e^{y/\beta}\PhiN(v)\right]
&=
\frac{1}{\beta} e^{y/\beta}\PhiN(v)
+
 e^{y/\beta} \phi(v) \frac{dv}{dy}
\notag\\
&=
\frac{1}{\beta} e^{y/\beta}\PhiN(v)
-
\frac{1}{\sigma} e^{y/\beta} \phi(v).
\label{eq:hprime-term2}
\end{align}
We now show that the Gaussian-density pieces cancel. Compute
\begin{align}
e^{-y/\beta} \phi(u)
&=
\frac{1}{\sqrt{2\pi}}
\exp\!\left(
-\frac{y}{\beta}
-\frac12 \left(\frac{y}{\sigma}-\frac{\sigma}{\beta}\right)^2
\right)
\notag\\
&=
\frac{1}{\sqrt{2\pi}}
\exp\!\left(
-\frac{y}{\beta}
-\frac{y^2}{2\sigma^2}
+\frac{y}{\beta}
-\frac{\sigma^2}{2\beta^2}
\right)
\notag\\
&=
\frac{1}{\sqrt{2\pi}}
\exp\!\left(-\frac{y^2}{2\sigma^2} - \frac{\sigma^2}{2\beta^2}\right).
\label{eq:phi-cancel-1}
\end{align}
Similarly,
\begin{align}
e^{y/\beta} \phi(v)
&=
\frac{1}{\sqrt{2\pi}}
\exp\!\left(
\frac{y}{\beta}
-\frac12 \left(-\frac{y}{\sigma}-\frac{\sigma}{\beta}\right)^2
\right)
\notag\\
&=
\frac{1}{\sqrt{2\pi}}
\exp\!\left(
\frac{y}{\beta}
-\frac{y^2}{2\sigma^2}
-\frac{y}{\beta}
-\frac{\sigma^2}{2\beta^2}
\right)
\notag\\
&=
\frac{1}{\sqrt{2\pi}}
\exp\!\left(-\frac{y^2}{2\sigma^2} - \frac{\sigma^2}{2\beta^2}\right).
\label{eq:phi-cancel-2}
\end{align}
Therefore the $\sigma^{-1}$ terms in \eqref{eq:hprime-term1} and \eqref{eq:hprime-term2} cancel. What remains is
\begin{equation}
h'(y;\beta,\sigma^2)
=
\frac{e^{\sigma^2/(2\beta^2)}}{2\beta^2}
\left[
 e^{y/\beta}\PhiN(v)
 -
 e^{-y/\beta}\PhiN(u)
\right].
\label{eq:hprime-final}
\end{equation}

\paragraph{Second derivative from an ODE identity.}
The centered Laplace density $\ell_\beta(y)=\frac{1}{2\beta}e^{-|y|/\beta}$ has Fourier transform
\[
\widehat{\ell_\beta}(\xi)=\frac{1}{1+\beta^2 \xi^2}.
\]
Therefore
\[
(1+\beta^2 \xi^2) \widehat{\ell_\beta}(\xi) = 1.
\]
Under the Fourier transform, multiplication by $1+\beta^2\xi^2$ corresponds to the operator $1-\beta^2 d^2/dy^2$. Hence, in the distributional sense,
\begin{equation}
\left(1-\beta^2 \frac{d^2}{dy^2}\right) \ell_\beta = \delta_0.
\label{eq:laplace-green}
\end{equation}
Now $h = \ell_\beta * \phi_{\sigma^2}$. Convolving \eqref{eq:laplace-green} with $\phi_{\sigma^2}$,
\[
\left(1-\beta^2 \frac{d^2}{dy^2}\right) h = \phi_{\sigma^2}.
\]
Therefore
\begin{equation}
h''(y;\beta,\sigma^2)
=
\frac{h(y;\beta,\sigma^2) - \phi_{\sigma^2}(y)}{\beta^2}.
\label{eq:hsecond-final}
\end{equation}

\section[Posterior of A0 given X0=x0 and exact density reduction]{Posterior of $A_0$ given $X_0=x_0$ and exact density reduction}
\label{sec:posterior-reduction}

From \eqref{eq:setup-ou},
\[
X_0 = e^{-t} A_0 + \sqrt{\Dt} Z_0.
\]
Hence
\[
X_0 \mid A_0=a_0 \sim \mathcal N(e^{-t} a_0, \Dt).
\]
The posterior density is proportional to prior times likelihood:
\begin{equation}
p(a_0\mid x_0)
\propto
\exp\!\left(
-\frac{(a_0-\mu_0)^2}{2\sigma_0^2}
-\frac{(x_0-e^{-t}a_0)^2}{2\Dt}
\right).
\label{eq:posterior-start}
\end{equation}

\subsection{Complete the square explicitly}

Expand both squares:
\begin{align}
-\frac{(a_0-\mu_0)^2}{2\sigma_0^2}
&=
-\frac{a_0^2}{2\sigma_0^2}
+\frac{\mu_0}{\sigma_0^2} a_0
-\frac{\mu_0^2}{2\sigma_0^2},
\label{eq:posterior-expand1}\\[0.25em]
-\frac{(x_0-e^{-t}a_0)^2}{2\Dt}
&=
-\frac{x_0^2}{2\Dt}
+\frac{e^{-t} x_0}{\Dt} a_0
-\frac{e^{-2t}}{2\Dt} a_0^2.
\label{eq:posterior-expand2}
\end{align}
Add \eqref{eq:posterior-expand1} and \eqref{eq:posterior-expand2}:
\begin{align}
&-\frac{(a_0-\mu_0)^2}{2\sigma_0^2} - \frac{(x_0-e^{-t}a_0)^2}{2\Dt}
\notag\\
&=
-\frac12 \left(\sigma_0^{-2} + e^{-2t}\Dt^{-1}\right) a_0^2
+
\left(\mu_0\sigma_0^{-2} + e^{-t} x_0 \Dt^{-1}\right) a_0
+
\text{constant in } a_0.
\label{eq:posterior-combined}
\end{align}
Define
\begin{equation}
v_t^{-1} := \sigma_0^{-2} + e^{-2t}\Dt^{-1},
\qquad
m_t(x_0) := v_t\left(\frac{\mu_0}{\sigma_0^2} + \frac{e^{-t} x_0}{\Dt}\right).
\label{eq:vt-mt-def}
\end{equation}
Then \eqref{eq:posterior-combined} becomes
\[
-\frac{1}{2v_t} a_0^2 + \frac{m_t(x_0)}{v_t} a_0 + \text{constant}.
\]
Complete the square:
\begin{align}
-\frac{1}{2v_t} a_0^2 + \frac{m_t}{v_t} a_0
&=
-\frac{1}{2v_t} \left(a_0^2 - 2m_t a_0\right)
\notag\\
&=
-\frac{1}{2v_t} \left((a_0-m_t)^2 - m_t^2\right)
\notag\\
&=
-\frac{(a_0-m_t)^2}{2v_t} + \frac{m_t^2}{2v_t}.
\label{eq:posterior-square}
\end{align}
Hence
\begin{equation}
A_0 \mid X_0=x_0 \sim \mathcal N(m_t(x_0),v_t).
\label{eq:posterior-final}
\end{equation}
Since $v_t$ depends only on $t$, not on $x_0$,
\begin{equation}
m_t'(x_0) = \frac{v_t e^{-t}}{\Dt}.
\label{eq:mt-prime}
\end{equation}

\paragraph{Marginal law of $X_0$.}
Directly from \eqref{eq:setup-ou},
\[
X_0 = e^{-t} A_0 + \sqrt{\Dt} Z_0
\]
is a sum of independent Gaussian variables, so
\begin{equation}
X_0 \sim \mathcal N(e^{-t}\mu_0, v_x),
\qquad
v_x := e^{-2t} \sigma_0^2 + \Dt.
\label{eq:x0-marginal-final}
\end{equation}

\subsection{Exact density reduction}

Start from the full joint integral over $a_0$ and $\varepsilon$:
\begin{align}
\pt(x_0,x_1)
&=
\int_{\R} \int_{\R}
\phi_{\sigma_0^2}(a_0-\mu_0)
\cdot
\frac{e^{-|\varepsilon|/b}}{2b}
\cdot
\phi_{\Dt}(x_0-e^{-t}a_0)
\cdot
\phi_{\Dt}(x_1-e^{-t}(\alpha a_0+\varepsilon))
\, d\varepsilon \, da_0.
\label{eq:pt-start}
\end{align}
Fix $a_0$ and define
\[
y = x_1 - e^{-t} \alpha a_0.
\]
Then
\[
x_1 - e^{-t}(\alpha a_0 + \varepsilon) = y - e^{-t} \varepsilon.
\]
The random variable $e^{-t}\varepsilon$ has Laplace law with scale $\beta_t=e^{-t}b$, so the $\varepsilon$-integral in \eqref{eq:pt-start} is exactly
\begin{equation}
h(y;\beta_t,\Dt) = h(x_1-e^{-t}\alpha a_0;\beta_t,\Dt).
\label{eq:innovation-integrated}
\end{equation}
Therefore \eqref{eq:pt-start} becomes
\begin{equation}
\pt(x_0,x_1)
=
\int_{\R}
\phi_{\sigma_0^2}(a_0-\mu_0)
\, \phi_{\Dt}(x_0-e^{-t}a_0)
\, h(x_1-c_t a_0;\beta_t,\Dt)
\, da_0,
\qquad c_t=e^{-t}\alpha.
\label{eq:pt-after-eps}
\end{equation}
By Bayes' rule,
\[
\phi_{\sigma_0^2}(a_0-\mu_0)\, \phi_{\Dt}(x_0-e^{-t}a_0)
=
p(x_0)\, p(a_0\mid x_0).
\]
We already computed both factors:
\[
p(x_0)=\phi_{v_x}(x_0-e^{-t}\mu_0),
\qquad
p(a_0\mid x_0)=\phi_{v_t}(a_0-m_t(x_0)).
\]
Substitute these into \eqref{eq:pt-after-eps}:
\begin{align}
\pt(x_0,x_1)
&=
\phi_{v_x}(x_0-e^{-t}\mu_0)
\int_{\R}
\phi_{v_t}(a_0-m_t(x_0))
\, h(x_1-c_t a_0;\beta_t,\Dt)
\, da_0
\notag\\
&=
\phi_{v_x}(x_0-e^{-t}\mu_0)
\; \E\!\left[h(x_1-c_t A_0;\beta_t,\Dt)\,\middle|\, X_0=x_0\right].
\label{eq:pt-exact}
\end{align}
This is the exact one-dimensional reduction used in the short note.

\section{Exact score and exact effective precision}
\label{sec:score-hessian}

Define, for $r=0,1,2$,
\begin{equation}
I_r(x_0,x_1;t)
:=
\E\!\left[h^{(r)}(x_1-c_t A_0;\beta_t,\Dt)\,\middle|\, X_0=x_0\right].
\label{eq:Ir-def-companion}
\end{equation}
Then \eqref{eq:pt-exact} reads
\begin{equation}
\pt(x_0,x_1) = \phi_{v_x}(x_0-e^{-t}\mu_0) I_0(x_0,x_1;t).
\label{eq:pt-I0}
\end{equation}

\subsection[Derivative rule for Ir with respect to x0]{Derivative rule for $I_r$ with respect to $x_0$}

Write the posterior variable as
\[
A_0 = m_t(x_0) + \sqrt{v_t} Z,
\qquad Z \sim \mathcal N(0,1).
\]
Then
\[
I_r(x_0,x_1;t)
=
\E\!\left[h^{(r)}\bigl(x_1-c_t(m_t(x_0)+\sqrt{v_t}Z)\bigr)\right].
\]
Because $v_t$ is independent of $x_0$, differentiation gives
\begin{align}
\partial_{x_0} I_r
&=
\E\!\left[
 h^{(r+1)}\bigl(x_1-c_t(m_t+\sqrt{v_t}Z)\bigr)
 \cdot
 (-c_t m_t'(x_0))
\right]
\notag\\
&=
-c_t m_t'(x_0) I_{r+1}(x_0,x_1;t).
\label{eq:Ir-x0-rule}
\end{align}
This identity is the key input for all score and Hessian calculations.

\subsection{Exact score formulas}

Take logarithms in \eqref{eq:pt-I0}:
\begin{equation}
\log \pt(x_0,x_1)
=
\log \phi_{v_x}(x_0-e^{-t}\mu_0)
+
\log I_0(x_0,x_1;t).
\label{eq:logp-split}
\end{equation}

\paragraph{Derivative with respect to $x_1$.}
Since the Gaussian prefactor depends only on $x_0$,
\begin{equation}
\partial_{x_1} \log \pt(x_0,x_1)
=
\frac{\partial_{x_1} I_0}{I_0}
=
\frac{I_1}{I_0}.
\label{eq:score-x1-final}
\end{equation}

\paragraph{Derivative with respect to $x_0$.}
First,
\[
\log \phi_{v_x}(x_0-e^{-t}\mu_0)
=
-\frac12 \log(2\pi v_x) - \frac{(x_0-e^{-t}\mu_0)^2}{2v_x}.
\]
Differentiate:
\begin{equation}
\partial_{x_0} \log \phi_{v_x}(x_0-e^{-t}\mu_0)
=
-\frac{x_0-e^{-t}\mu_0}{v_x}.
\label{eq:gaussian-prefactor-deriv}
\end{equation}
Also, by \eqref{eq:Ir-x0-rule} with $r=0$,
\[
\partial_{x_0} I_0 = -c_t m_t'(x_0) I_1.
\]
Hence
\begin{equation}
\partial_{x_0} \log \pt(x_0,x_1)
=
-\frac{x_0-e^{-t}\mu_0}{v_x}
-c_t m_t'(x_0) \frac{I_1}{I_0}.
\label{eq:score-x0-final}
\end{equation}
Equations \eqref{eq:score-x1-final} and \eqref{eq:score-x0-final} are the exact score formulas used in the short note.

\subsection{Exact Hessian formulas}

Define the effective precision
\begin{equation}
\Ht(x) := -\nabla_x^2 \log \pt(x).
\label{eq:H-def}
\end{equation}
We now compute its entries explicitly.

\paragraph{The $(1,1)$ entry.}
Differentiate \eqref{eq:score-x1-final} with respect to $x_1$:
\begin{align}
\partial_{x_1 x_1} \log \pt
&=
\partial_{x_1}\left(\frac{I_1}{I_0}\right)
\notag\\
&=
\frac{(\partial_{x_1} I_1) I_0 - I_1 (\partial_{x_1} I_0)}{I_0^2}
\notag\\
&=
\frac{I_2 I_0 - I_1^2}{I_0^2}.
\label{eq:L11}
\end{align}
Therefore
\begin{equation}
(\Ht(x))_{11}
=
-\partial_{x_1 x_1} \log \pt(x)
=
\frac{I_1^2 - I_0 I_2}{I_0^2}.
\label{eq:H11-final}
\end{equation}

\paragraph{The mixed entry.}
Differentiate \eqref{eq:score-x1-final} with respect to $x_0$ and use \eqref{eq:Ir-x0-rule} twice:
\begin{align}
\partial_{x_0 x_1} \log \pt
&=
\partial_{x_0}\left(\frac{I_1}{I_0}\right)
\notag\\
&=
\frac{(\partial_{x_0} I_1) I_0 - I_1 (\partial_{x_0} I_0)}{I_0^2}
\notag\\
&=
\frac{\bigl(-c_t m_t' I_2\bigr) I_0 - I_1 \bigl(-c_t m_t' I_1\bigr)}{I_0^2}
\notag\\
&=
-c_t m_t'(x_0) \frac{I_2 I_0 - I_1^2}{I_0^2}
\notag\\
&=
 c_t m_t'(x_0) (\Ht(x))_{11}.
\label{eq:L01}
\end{align}
Therefore
\begin{equation}
(\Ht(x))_{01}
=
-\partial_{x_0 x_1} \log \pt(x)
=
-c_t m_t'(x_0) (\Ht(x))_{11}.
\label{eq:H01-final}
\end{equation}

\paragraph{The $(0,0)$ entry.}
Differentiate \eqref{eq:score-x0-final} with respect to $x_0$:
\begin{align}
\partial_{x_0 x_0} \log \pt
&=
-\frac{1}{v_x}
-c_t m_t'(x_0) \partial_{x_0}\left(\frac{I_1}{I_0}\right)
\label{eq:L00-step1}
\end{align}
because $m_t'(x_0)$ is constant in $x_0$. Using \eqref{eq:L01},
\[
\partial_{x_0}\left(\frac{I_1}{I_0}\right) = c_t m_t'(x_0) (\Ht(x))_{11}.
\]
Substitute into \eqref{eq:L00-step1}:
\begin{equation}
\partial_{x_0 x_0} \log \pt
=
-\frac{1}{v_x} - c_t^2 \bigl(m_t'(x_0)\bigr)^2 (\Ht(x))_{11}.
\label{eq:L00-final}
\end{equation}
Therefore
\begin{equation}
(\Ht(x))_{00}
=
\frac{1}{v_x} + c_t^2 \bigl(m_t'(x_0)\bigr)^2 (\Ht(x))_{11}.
\label{eq:H00-final}
\end{equation}
This completes the derivation of the exact curvature matrix.

\subsection{Why constant Hessian would force Gaussianity}

The next statement is the clean logical bridge from the formulas above to the presentation claim.

\begin{proposition}
Let $p$ be a strictly positive $C^2$ density on $\R^n$. If $-\nabla^2 \log p(x) \equiv H$ is constant, then $p$ is Gaussian.
\end{proposition}

\begin{proof}
If $-\nabla^2 \log p(x) \equiv H$, then
\[
\nabla^2 \log p(x) = -H
\qquad \text{for all } x.
\]
Integrating once, there exists a vector $b\in\R^n$ such that
\[
\nabla \log p(x) = -H x + b.
\]
Now integrate along the straight segment $s \mapsto s x$ from $0$ to $x$:
\begin{align}
\log p(x) - \log p(0)
&=
\int_0^1 \frac{d}{ds} \log p(sx)\, ds
\notag\\
&=
\int_0^1 \langle \nabla \log p(sx), x \rangle \, ds
\notag\\
&=
\int_0^1 \langle -H(sx)+b, x \rangle \, ds
\notag\\
&=
-\int_0^1 s \, x^\top H x \, ds + \int_0^1 b^\top x \, ds
\notag\\
&=
-\frac12 x^\top H x + b^\top x.
\label{eq:gaussianity-proof}
\end{align}
Hence
\[
p(x) = C \exp\!\left(-\frac12 x^\top H x + b^\top x\right)
\]
for some constant $C>0$. Completing the square shows that this is a Gaussian density.\qedhere
\end{proof}

Apply the proposition to $p=\pt$. Since the characteristic function in \eqref{eq:cf-final} is not Gaussian for finite $t$, the Hessian cannot be constant. Therefore the exact formulas \eqref{eq:H11-final}--\eqref{eq:H00-final} are not merely decorative: they encode a genuinely non-Gaussian effect.

\section[What can be stated exactly for K=2]{What can be stated exactly for $K=2$}
\label{sec:k2-statement}

Let
\[
A_0 \sim \mathcal N(\mu_0,\sigma_0^2),
\qquad
A_1 = \alpha A_0 + \varepsilon_1,
\qquad
A_2 = \alpha A_1 + \varepsilon_2,
\]
with $\varepsilon_1,\varepsilon_2 \stackrel{\mathrm{iid}}{\sim} \mathrm{Laplace}(0,b)$, followed by the same OU noising on each coordinate. Then the exact noisy density is
\begin{align}
\pt(x_0,x_1,x_2)
&=
\int_{\R} \phi_{\sigma_0^2}(a_0-\mu_0) \phi_{\Dt}(x_0-e^{-t}a_0)
\; J_t(a_0;x_1,x_2)\, da_0,
\label{eq:k2-outer}
\end{align}
where
\begin{align}
J_t(a_0;x_1,x_2)
&=
\int_{\R} \frac{e^{-|a_1-\alpha a_0|/b}}{2b}
\, \phi_{\Dt}(x_1-e^{-t}a_1)
\, h(x_2-e^{-t}\alpha a_1;\beta_t,\Dt)
\, da_1.
\label{eq:k2-inner}
\end{align}
Formulas \eqref{eq:k2-outer}--\eqref{eq:k2-inner} are exact. However, this is still a nested two-layer integral representation. The current project does \emph{not} claim an elementary closed form comparable to the $K=1$ formula \eqref{eq:pt-exact}. That stronger claim would require additional work.

\section{Numerical realization used in the figures}
\label{sec:numerics}

The final figures do not use Monte Carlo. They evaluate the exact formulas above by Gauss-Hermite quadrature.

If $A\sim\mathcal N(m,v)$ and $f$ is integrable, then
\begin{align}
\E[f(A)]
&=
\int_{\R} f(a) \frac{1}{\sqrt{2\pi v}} e^{-(a-m)^2/(2v)} da.
\label{eq:gh-start}
\end{align}
Set
\[
a = m + \sqrt{2v}\, z,
\qquad da = \sqrt{2v} \, dz.
\]
Then
\begin{align}
\E[f(A)]
&=
\int_{\R} f(m+\sqrt{2v}z) \frac{1}{\sqrt{2\pi v}} e^{-z^2} \sqrt{2v} \, dz
\notag\\
&=
\frac{1}{\sqrt{\pi}} \int_{\R} f(m+\sqrt{2v}z) e^{-z^2} dz.
\label{eq:gh-transform}
\end{align}
If $(z_i,w_i)_{i=1}^M$ are Hermite nodes and weights for $\int e^{-z^2} g(z) dz$, then
\begin{equation}
\E[f(A)] \approx \frac{1}{\sqrt{\pi}} \sum_{i=1}^M w_i f(m+\sqrt{2v}z_i).
\label{eq:gh-quadrature}
\end{equation}
This is exactly the formula used in the code to compute $I_0$, $I_1$, and $I_2$.

The code also includes independent checks:
\begin{enumerate}[leftmargin=1.45em,itemsep=0.25em]
\item direct inversion checks for the Gaussian finite-chain precision and for the deterministic control;
\item a direct convolution integral confirming the corrected normal-Laplace formula;
\item an independent double integral confirming the $K=1$ reduction formula;
\item finite-difference checks confirming the exact score and Hessian formulas.
\end{enumerate}
The script \texttt{code/run\_checks.py} writes the resulting errors to \texttt{code/check\_report.txt}.

\section{Audit summary}

The final short note keeps only the claims that are either exact or clearly marked.
\begin{enumerate}[leftmargin=1.45em,itemsep=0.28em]
\item \textbf{Kept exactly:}
\begin{itemize}[leftmargin=1.2em]
\item the Gaussian OU identities $\Sigma_t=e^{-2t}\Sigma_0+\Dt I$ and $Q_t=\Sigma_t^{-1}$;
\item the small-time expansion $Q_t=e^{2t}\sum_{m\ge 0}(-c_t)^m Q_0^{m+1}$ under the convergence condition $c_t\|Q_0\|<1$;
\item the large-time expansion $Q_t=I-e^{-2t}(\Sigma_0-I)+O(e^{-4t})$;
\item the deterministic-control inverse formula \eqref{eq:det-precision};
\item the exact Laplace $K=1$ characteristic function, density reduction, score, and effective-precision formulas.
\end{itemize}
\item \textbf{Removed or weakened:}
\begin{itemize}[leftmargin=1.2em]
\item the inaccurate boundary form of the finite-chain Gaussian precision;
\item the claim that the original locality observables ``fully characterize locality'';
\item the interpretation that deterministic dense precision and stochastic Markov fill-in are the same mechanism;
\item the old $K=2$ elementary closed-form claim.
\end{itemize}
\item \textbf{Main message now supported rigorously:}
\[
\begin{aligned}
\text{Gaussian benchmark} &\Longrightarrow -\nabla^2 \log p_t \equiv Q_t,\\
\text{Laplace benchmark} &\Longrightarrow -\nabla^2 \log p_t(x) = \Ht(x) \text{ depends on } x.
\end{aligned}
\]
This is the central statement that the final short presentation is built to communicate.
\end{enumerate}

\end{document}
